%  =============================================================================
%  Datei-Hinweise & Konfiguration: siehe config.md
%  -> Abschnitt "kapitel/04_schluss.tex"
%  =============================================================================

\chapter{Schluss}
\label{cha:schluss}

Nachdem der Handelsagent entworfen, implementiert und evaluiert wurde, fasst dieses
Kapitel die zentralen Erkenntnisse mit Blick auf die Zielsetzung zusammen, diskutiert
Stärken und Schwächen der gewählten Lösung und gibt einen Ausblick auf
weiterführende Arbeiten.

\section{Fazit}
\label{sec:fazit}

Ziel dieser Arbeit war der Aufbau eines Live-Trading-Agenten auf Basis des
Agenten-Frameworks OpenClaw, dessen deterministische \glsxtrshort{sma}-Strategie
durch ein Gremium großer Sprachmodelle geprüft wird, bevor ein Mehrheitsentscheid
die Orderausführung freigibt (\cref{sec:zielsetzung}). Dieses Ziel wurde erreicht:
Der vollständige Handelszyklus -- von der kontinuierlichen, gegen eine
Referenzreihe abgesicherten Marktdatenerfassung über die regelbasierte
Regimeentscheidung und das Multi-Modell-Voting bis zur Orderaufgabe über die
Schnittstelle einer regulierten Bank -- ist umgesetzt, im Trockenlauf erprobt und
in \cref{cha:umsetzung} vollständig dokumentiert. Die theoretischen Grundlagen aus
\cref{cha:grundlagen} bildeten dafür den methodischen Rahmen: Das Periodensystem
der Künstlichen Intelligenz erlaubte es, die bewusste Begrenzung des
\glsxtrshort{ki}-Einsatzes nicht als Verzicht, sondern als gestalterische Aussage
zu formulieren -- die Elemente der Kursprognose und der selbstlernenden
Fortschreibung bleiben unbesetzt, das Element der prozesssteuernden Kontrolle
bildet das Zentrum.

Mit Blick auf die erste Leitfrage zeigt die Fallstudie, dass der autonome,
nachvollziehbare Betrieb mit vertretbarem Aufwand möglich ist -- der wesentliche
Aufwand liegt dabei nicht im Handelsweg, sondern in seiner Absicherung: im
zweiphasigen Ausführungsjournal gegen Doppelorders, im Orderbuchabgleich nach
Kommunikationsfehlern, in den Kapitalobergrenzen sowie im jederzeit wirksamen
Not-Aus (\cref{sec:umsetzung-implementierung}). Die zweite Leitfrage nach dem
messbaren Beitrag des Mehrheitsvotums bleibt zum Abgabezeitpunkt bewusst offen;
die dafür geschaffene Kalibrierung misst den Veto-Beitrag der Modelle fortlaufend
gegen die tatsächliche Marktentwicklung und benötigt eine längere Betriebsdauer,
bevor valide Aussagen möglich sind (\cref{sec:umsetzung-evaluation-design}).

In Kennzahlen verdichtet stellt sich das Ergebnis wie folgt dar: Der Backtest
über zehn Jahre stellt dem täglichen Round-Trip der Demonstrationsstrategie
(rund $-29.700$\,€, davon etwa 98\,\% Transaktionskosten) das reine Regimefolgen
(rund $+96\,\%$ vor Steuern bei neun Positionswechseln) und das passive Halten
(rund $+248\,\%$ vor Steuern) gegenüber; die Walk-Forward-Prüfung des
Prognosemodells über
23.803 Handelstage ergibt keine Prognosekraft über die Basisrate hinaus bei
zugleich guter Kalibrierung; die real gemessenen Fixkosten betragen 7,80\,€ je
Round-Trip, und die Antwortlatenzen des Modellgremiums liegen im eingeschwungenen
Betrieb unter einer Sekunde je Votum
(\cref{sec:umsetzung-evaluation-ergebnisse}). Ökonomisch fällt die Antwort damit
eindeutig aus: Die täglich handelnde Demonstrationsstrategie ist nach realen
Kosten strukturell defizitär; der Agent ist ein Erkenntnis-, kein
Renditeinstrument, und die Differenz zum reinen Regimefolgen beziffert den Preis
der täglich demonstrierten Autonomie. Methodisch hat sich dabei das iterative
Vorgehen bewährt, jede Ausbaustufe unmittelbar gegen reale Betriebsdaten zu
prüfen: Mehrere der wirksamsten Absicherungen -- von der Unterscheidung legitimer
Datenstille von echten Verbindungsverlusten bis zur datumskorrekten
Zweitquellenprüfung -- entstanden nicht am Reißbrett, sondern aus Befunden des
laufenden Betriebs und wurden anschließend durch Regressionstests dauerhaft
abgesichert.

\section{Diskussion}
\label{sec:diskussion}

Den Stärken der gewählten Lösung stehen benennbare Schwächen gegenüber; beide
werden hier bewusst gebündelt, um die Zielerreichung ehrlich einzuordnen. Als
tragfähig erwiesen hat sich erstens das Prinzip der begrenzten Diskretion: Kein
beobachteter Fehlerfall -- weder Schemaverletzungen einzelner Modelle noch
Zeitüberschreitungen oder Datenlücken -- konnte eine unbeabsichtigte Order
auslösen, weil jede Unklarheit konservativ als Veto oder Handelsstopp endet.
Zweitens hat sich die Entscheidung ausgezahlt, die Evaluation auf gemessene statt
angenommene Größen zu stützen: Reale Kostenausweise, die laufende
Kursgegenprüfung und die Walk-Forward-Prüfung machen die Ergebnisse belastbar, und der ehrliche
Negativbefund zur Prognosekraft trägt die zentrale Architekturentscheidung
empirisch, statt sie zu beschädigen. Drittens schafft die durchgängige
Journalisierung zusammen mit dem öffentlichen Dashboard eine Transparenz, die
jede einzelne Entscheidung samt Modellvoten nachvollziehbar und den laufenden
Betrieb öffentlich prüfbar macht.

Dem stehen Schwächen gegenüber. Der erhoffte Sicherheitsgewinn des
Modellgremiums ist zu Betriebsbeginn empirisch unbelegt, und korrelierte
Fehlurteile bleiben trotz Anbieterdiversität möglich, da alle Modelle denselben
Entscheidungskontext erhalten; die Kalibrierung kann dies erst nach längerer
Betriebsdauer klären. Die tägliche Handelsfrequenz ist nach Kosten strukturell
defizitär und nur als Demonstrationsvehikel zu rechtfertigen. Die Autonomie
bleibt an regulatorische und betriebliche Grenzen gebunden -- die tägliche
photoTAN-Freigabe, den bewusst manuell gehaltenen Ausstiegspfad bei
Regimewechseln und eine Datenfrische von 60 bis 90 Sekunden am
Entscheidungspunkt. Schließlich stützt sich die Kursdatenerfassung auf eine
inoffizielle Quelle ohne Verfügbarkeitsgarantie, deren Ausfall der Betrieb zwar
erkennen, aber nicht kompensieren kann. Aus dieser Gegenüberstellung folgt das
Verbesserungspotenzial, das der nachstehende Ausblick in konkrete
Weiterentwicklungen übersetzt.

\section{Ausblick}
\label{sec:ausblick}

Aus der Bearbeitung haben sich die folgenden weiterführenden Ansatzpunkte ergeben,
die Gegenstand künftiger Arbeiten sein können. Als Erstes zu nennen ist der
unmittelbar bevorstehende Übergang vom Trockenlauf in den Echtgeldbetrieb, der die
bislang bis zur Ordervalidierung demonstrierte Ausführungsseite unter realen
Bedingungen prüft und die Datenbasis der Voting-Kalibrierung kontinuierlich
erweitert. Als weiterer Punkt kann der Soll-Positions-Abgleich umgesetzt werden,
der den Ausstiegspfad bei Regimewechseln vollständig deterministisch macht und
damit die letzte bewusst manuell gehaltene Zwischenlösung ablöst; darauf aufbauend
erlaubt der Parallelbetrieb eines zweiten Depots den direkten Vergleich der täglich
handelnden Kreuzungsstrategie mit der haltenden Regimestrategie unter identischen
Bedingungen. Des Weiteren bietet sich die Weiterentwicklung des Modellgremiums hin
zu komplementären Prüfrollen an, um der Gefahr korrelierter Fehlurteile zu
begegnen; die laufende Kalibrierung liefert hierfür die datenbasierte
Entscheidungsgrundlage, welche Modelle das Gremium tatsächlich stärken.
Abschließend wäre die Migration auf einen Mac~mini mit lokal betriebenen Modellen
anzustreben, die die laufenden Kosten weitgehend eliminiert und die Datenhoheit
vervollständigt -- die Architektur ist darauf vorbereitet, der Wechsel bleibt eine
Konfigurationsentscheidung.

Damit liegt mit dieser Arbeit eine nachvollziehbar dokumentierte, übertragbare
Referenzarchitektur für sicherheitsorientiert begrenzten \glsxtrshort{ki}-Einsatz
im autonomen Handel vor, die als Grundlage für die weiterführende Untersuchung der
Frage dient, unter welchen Bedingungen probabilistische Sprachmodelle als
kontrollierte Prüfinstanzen deterministischer Systeme einen messbaren
Sicherheitsgewinn stiften.

%  =============================================================================
%  Ende
%  =============================================================================
